\newif\ifRELEASE
%\RELEASEfalse
\RELEASEtrue

\documentclass[10pt]{beamer}

\usepackage{amsmath,amssymb,enumerate,calc,color,ifthen,capt-of,booktabs,graphicx,listings,palatino,amsbsy,subfigure}

\usefonttheme{serif}
 
\definecolor{light-gray}{gray}{0.95}
\lstset{basicstyle=\footnotesize\ttfamily,
        %numbers=left,
        %backgroundcolor=\color{light-gray},
        frame=single, 
        %rulesepcolor=\color{red},
        keywordstyle=\color[rgb]{0,0,1},
        commentstyle=\color[rgb]{0.133,0.545,0.133},
        stringstyle=\color[rgb]{0.627,0.126,0.941},
        captionpos=b,
        title=\lstname,
        showstringspaces=false
       }


%---TITLE AND AUTHOR INFORMATION
\title % (optional, use only with long paper titles)
[OpenMP]{Introduction to OpenMP}
\author[Roberts]{Prof.~Jeremy Roberts}
% \institute[22.213] % (optional, but mostly needed)
%  {}
% - Keep it simple, no one is interested in your street address.
\date
%[CFP 2003] % (optional, should be abbreviation of conference name)
{11/18/2015}


\begin{document}

% TITLE PAGE
\begin{frame}[plain]
  \titlepage
\end{frame}

% What are the realities of full core, time-dependent simulations?

%------------------------------------------------------------------------------%
\begin{frame}{Resources}

\begin{itemize}
 \item Chapman Jost, and Van Der Pas, {\it Using OpenMP}, MIT Press (2008)
 \item ``Parallel Programming for Multicore Machines Using OpenMP and MPI''
       as taught by Dr. Evangelinos, MIT OpenCourseWare (2010)
       \begin{itemize}
         \item \url{ocw.mit.edu/courses/earth-atmospheric-and-planetary-sciences/12-950-parallel-programming-for-multicore-machines-using-openmp-and-mpi-january-iap-2010/}
       \end{itemize}
 \item OpenMP cheat sheets for C/C++ and Fortran:
       \begin{itemize}
         \item \url{openmp.org/mp-documents/OpenMP-4.0-C.pdf}
         \item \url{openmp.org/mp-documents/OpenMP-4.0-Fortran.pdf}
       \end{itemize}
\end{itemize}
\vfill
For next time: \textcolor{red}{\tt conda install mpi4py}

\end{frame}

%------------------------------------------------------------------------------%
\begin{frame}{OpenMP Programming Model}

\begin{itemize}
 \item OpenMP defines an application programmer interface (API) for 
       shared-memory, parallel programming based on the ``fork-join'' model
 \item Threads operate within the same address space (i.e., shared memory)
 \item Coarse- and fine-grain parallelism is supported
 \item Threads read/write shared variables (but watch out 
       for \textcolor{red}{race conditions}!) 
       and private variables
\item  Long-term development and support for C/C++ and Fortran (i.e.,
       OpenMP will be around for a long time)
\end{itemize}

\end{frame}

%------------------------------------------------------------------------------%
\begin{frame}{Hello World}

\begin{columns}[c]
  \begin{column}{0.5\textwidth}
    \lstinputlisting[language=C++]{code/hello.cc}
  \end{column}
  \begin{column}{0.5\textwidth}
    \lstinputlisting[language=Fortran]{code/hello.f90}
  \end{column}
\end{columns}
\vfill 
Watch for your first \textcolor{red}{race condition}
in the Fortran example!

\end{frame}

%------------------------------------------------------------------------------%
\begin{frame}{Compiling Your First OpenMP Program}

For C++, use (in the command line)
\begin{equation*}
 \overbrace{\tt g\!+\!+}^{\text{compiler}}~
   \textcolor{red}{\tt{-fopenmp}}~
   \underbrace{\tt{hello.cc}}_{\text{file to compile}}~
     \overbrace{\tt{-o}}^{\text{output as}}~
       \underbrace{\tt{hello\_c}}_{\text{this executable}} 
\end{equation*}
\vfill 

For Fortran, use
\begin{equation*}
 \overbrace{\tt gfortran}^{\text{compiler}}~
   \textcolor{red}{\tt{-fopenmp}}~
   \underbrace{\tt{hello.f90}}_{\text{file to compile}}~
     \overbrace{\tt{-o}}^{\text{output as}}~
       \underbrace{\tt{hello\_f}}_{\text{this executable}} 
\end{equation*}

\vfill 
The \textcolor{red}{\tt{-fopenmp}} automatically enables threading
in your code, though you need \textcolor{red}{\tt \#include <omp.h>} or
\textcolor{red}{\tt use omp\_lib} for use 
of OpenMP functions.  (Note also existence of \textcolor{red}{\tt \_OPENMP}
macro.)
\end{frame}

%------------------------------------------------------------------------------%
\begin{frame}{Simplest Real-World Use: Sum an Array}

\begin{columns}[c]
  \begin{column}{0.5\textwidth}
    \lstinputlisting[language=C++]{code/sum_serial.cc}
  \end{column}
  \begin{column}{0.5\textwidth}
    \lstinputlisting[language=Fortran]{code/sum_serial.f90}
  \end{column}
\end{columns}

\end{frame}

%------------------------------------------------------------------------------%
\begin{frame}{A N\"{a}ive Attempt...}

\begin{columns}[c]
  \begin{column}{0.5\textwidth}
    \lstinputlisting[language=C++]{code/sum_thread_1.cc}
  \end{column}
  \begin{column}{0.5\textwidth}
    \lstinputlisting[language=Fortran]{code/sum_thread_1.f90}
  \end{column}
\end{columns}

\end{frame}

%------------------------------------------------------------------------------%
\begin{frame}{Correct, But Bad Performance!}

\begin{columns}[c]
  \begin{column}{0.5\textwidth}
    \lstinputlisting[language=C++]{code/sum_thread_2.cc}
  \end{column}
  \begin{column}{0.5\textwidth}
    \lstinputlisting[language=Fortran]{code/sum_thread_2.f90}
  \end{column}
\end{columns}

\end{frame}

%------------------------------------------------------------------------------%
\begin{frame}{Correct, With Better Performance!}


\begin{columns}[c]
  \begin{column}{0.5\textwidth}
    \lstinputlisting[language=C++,basicstyle=\tiny\ttfamily]{code/sum_thread_3.cc}
  \end{column}
  \begin{column}{0.5\textwidth}
    \lstinputlisting[language=Fortran,basicstyle=\tiny\ttfamily]{code/sum_thread_3.f90}
  \end{column}
\end{columns}

\end{frame}

%------------------------------------------------------------------------------%
\begin{frame}{Correct, With Better Performance!}


\lstinputlisting[bash,basicstyle=\tiny\ttfamily]{code/results.sh}


\end{frame}

%------------------------------------------------------------------------------%
\begin{frame}{Useful OpenMP Directives}

\lstinputlisting[language=C++]{code/useful_directives.cc}

\end{frame}

%------------------------------------------------------------------------------%
\begin{frame}{Useful OpenMP Functions}

\lstinputlisting[language=C++]{code/useful_functions.cc}

\end{frame}

%------------------------------------------------------------------------------%
\begin{frame}{Useful OpenMP Environment Variables}

\lstinputlisting[language=bash]{code/useful_env_vars.sh}

\end{frame}

\end{document}
